\section{总结与延伸}

\subsection{三大核心原则}

Anthropic 在文章结尾提炼了构建 Agent 的三大核心原则：

\begin{importantbox}{构建 Agent 的三大核心原则}
\begin{enumerate}
    \item \textbf{保持设计简洁}（Maintain simplicity）：Agent 的架构应该尽可能简单，只在证明有效的情况下增加复杂度。
    \item \textbf{优先透明性}（Prioritize transparency）：明确展示 Agent 的规划步骤，让用户和开发者能理解 Agent 的决策过程。
    \item \textbf{精心打造 ACI}（Carefully craft ACI）：通过完善的工具文档和测试，确保 Agent 能正确、高效地与外部世界交互。
\end{enumerate}
\end{importantbox}

\subsection{从本文可以学到什么}

这篇文章的价值不仅在于具体的设计模式，更在于其传达的\textbf{工程哲学}：

\begin{enumerate}
    \item \textbf{从简到繁}：不要被「Agent」的概念所迷惑，先用最简单的方案尝试，再逐步升级。
    \item \textbf{模式思维}：五种 Workflow 模式和 Agent 模式提供了一个实用的思考框架，帮助开发者在面对新任务时快速匹配合适的架构。
    \item \textbf{工具即界面}：Agent 的能力上限很大程度上取决于工具设计的质量，而非模型本身或 prompt 的精巧程度。
    \item \textbf{人机协作}：即使是最自主的 Agent，也应该设计人类监督和干预的机制。
\end{enumerate}

\subsection{拓展阅读}

\begin{itemize}
    \item Anthropic 官方 Cookbook：\url{https://github.com/anthropics/anthropic-cookbook} ——包含各种 agentic 模式的代码实现示例。
    \item Model Context Protocol (MCP)：\url{https://modelcontextprotocol.io/} ——Agent 工具生态的统一标准。
    \item SWE-bench：\url{https://www.swebench.com/} ——评估编程 Agent 能力的标准基准。
    \item Claude Agent SDK：\url{https://github.com/anthropics/claude-agent-sdk} ——Anthropic 官方 Agent 开发工具包。
\end{itemize}

%% --- Fin del contenido --- %%

\end{document}
